\documentclass[11pt]{article}

\usepackage[margin=1in]{geometry}
\usepackage{amsmath,amssymb,amsthm}
\usepackage{enumitem}
\usepackage[hidelinks]{hyperref}

\newtheorem{theorem}{Theorem}
\newtheorem{proposition}[theorem]{Proposition}
\newtheorem{corollary}[theorem]{Corollary}
\newtheorem{lemma}[theorem]{Lemma}
\theoremstyle{definition}
\newtheorem{definition}[theorem]{Definition}
\theoremstyle{remark}
\newtheorem{remark}[theorem]{Remark}

\newcommand{\Z}{\mathbb Z}
\newcommand{\F}{\mathbb F}
\newcommand{\D}{\mathbb Z[1/2]}
\newcommand{\Cl}{\operatorname{Cl}}
\newcommand{\ShUg}{\operatorname{ShUg}}
\newcommand{\Tor}{\operatorname{Tor}}
\newcommand{\taglabel}[1]{\textnormal{\textbf{[#1]}}}
\newcommand{\PROVED}{\taglabel{PROVED}}
\newcommand{\SOURCE}{\taglabel{SOURCE-PINNED}}
\newcommand{\BOUNDARY}{\taglabel{BOUNDARY}}

\title{The divisibility obstruction to a game-native Clifford deformation}
\author{Codex}
\date{2026-08-09}

\begin{document}
\maketitle

\begin{abstract}
The exterior algebra of the short-game group admits many quadratic
deformations after one first chooses a small subgroup and then supplies a
quadratic table by hand.  That local freedom does not extend to a quadratic
datum belonging to game values themselves.  The reason is the abstract group
structure of short games: it is $2$-divisible and all of its torsion is
$2$-primary.  We prove that in any coefficient-faithful Clifford realization
of such a group, every torsion game has square zero and is polar-orthogonal to
every game.  Equivalently, every globally defined, or subgroup-inclusion
coherent, quadratic datum factors through the torsion-free quotient.  This
settles the \texttt{tisn} problem negatively under an exact naturality
condition and explains why the nonzero Gold forms on nimbers cannot extend to
general partizan games.
\end{abstract}

\section{The question and the naturality contract}

The implemented object \path{GameExterior} starts with finitely many short
games, takes their $\Z$-span $M$, and forms $\Lambda_{\Z}M$.  The checked
\path{GameClifford} surface accepts integer values
\[
 e_x^2=Q(x),\qquad e_xe_y+e_ye_x=B(x,y)
\]
only after every known relation in $M$ is null and polar-radical.  Its tables
are intentionally supplied by the caller.  This leaves a mathematical
question: can the game structure itself provide a nonzero torsion square,
rather than merely allowing a table on one selected presentation?

The phrase ``game-native'' needs a naturality condition.  Without one, the
question has the tautological answer ``yes'': on $\langle *\rangle\cong\Z/2$
one may choose characteristic-two coefficients and declare $e_*^2=1$.
Nothing in that declaration distinguishes $*$ from an abstract generator of
$\Z/2$, and the value changes when the ambient subgroup changes.

\begin{definition}[ambient coherence]\label{def:coherent}
A Clifford deformation of short games is \emph{ambient-coherent} if it is
either a single coefficient-faithful datum on the full additive group
$\ShUg$, or, sufficiently for finite calculations, a directed family of data
on finitely generated subgroups $M\leq\ShUg$ such that every inclusion
$M\subseteq N$ induces injective coefficient and algebra maps preserving the
grade-one generators, their squares, and their polar anticommutators.
\end{definition}

This is the weakest contract under which $Q(x)$ is attached to the game value
$x$ rather than to the accident that one stopped at a root-incomplete
subgroup.  It permits the coefficient ring and algebra to grow with the
subgroup.  Injectivity is essential: without it an inclusion may ``preserve''
a square only by killing the coefficient that recorded it.

\section{The external game-group input}

Let $\D=\Z[1/2]$ be the additive group of dyadic rationals.

\begin{theorem}[Moews]\label{thm:moews}
\SOURCE{} The additive group of short partizan games has abstract group
structure
\[
 \ShUg\cong\bigoplus_{\mathbb N}\D\;\oplus\;
              \bigoplus_{\mathbb N}\D/\Z.
\]
Moreover, every finite-order short game has order a power of two.
\end{theorem}

The order statement is Theorem~8 and the decomposition is Theorem~9 of
Moews~\cite{moews}.  Two consequences, and only these consequences, enter the
new proof:
\begin{enumerate}[label=(\roman*),leftmargin=2.5em]
\item multiplication by $2^k$ on $\ShUg$ is surjective for every $k$;
\item if $t$ is torsion, then $2^k t=0$ for some $k$.
\end{enumerate}
Notice that $\ShUg$ is not divisible by every integer: the $\D$ summands are
not $3$-divisible.  Full divisibility is neither claimed nor needed.

\section{The algebraic root lemma}

The obstruction occurs before any classification of quadratic maps.  It is a
statement about an additive grade-one realization in an arbitrary associative
ring.

\begin{lemma}[root collapse]\label{lem:root}
\PROVED{} Let $A$ be an abelian group, $C$ an associative ring, and
$i:A\to(C,+)$ an additive map.  Suppose $n t=0$ and choose $y,z\in A$ with
\[
 ny=t,\qquad nz=x.
\]
Then
\[
 i(t)^2=0,
 \qquad
 i(t)i(x)+i(x)i(t)=0.
\]
\end{lemma}

\begin{proof}
The relation $nt=0$ and the identity $ny=t$ give
\[
 n^2y=n(ny)=nt=0,
\]
hence $n^2i(y)=i(n^2y)=0$.  Since integer scalars are central in $C$,
\[
\begin{aligned}
 i(t)^2
   &=(n i(y))(n i(y))
     =n^2 i(y)^2
     =(n^2i(y))i(y)=0,\\
 i(t)i(x)+i(x)i(t)
   &=(n i(y))(n i(z))+(n i(z))(n i(y))\\
   &=n^2\bigl(i(y)i(z)+i(z)i(y)\bigr)\\
   &=(n^2i(y))i(z)+i(z)(n^2i(y))=0.
\end{aligned}
\]
No commutativity of $C$, quadratic homogeneity axiom, or assumption on its
characteristic is used.
\end{proof}

\begin{remark}[formal proof]
\PROVED{} The kernel-checked file is \texttt{GameExterior.lean}.  Its theorem takes the roots
$y,z$ explicitly, proves both ring identities, and then proves the
coefficient-valued corollaries below from injectivity of the algebra map.  The
file contains no \path{sorry}, custom axiom, or finite computation.  The
Moews theorem remains an explicit external input rather than being encoded as
an axiom.
\end{remark}

\section{Global torsion-radical collapse}

Let $R$ be a commutative ring and $C$ an $R$-algebra.  A
coefficient-faithful Clifford datum consists of an additive grade-one map
$i:A\to C$, functions $Q:A\to R$ and $B:A\times A\to R$, relations
\[
 i(x)^2=\iota(Q(x)),\qquad
 i(x)i(y)+i(y)i(x)=\iota(B(x,y)),
\]
and an injective coefficient map $\iota:R\hookrightarrow C$.

\begin{theorem}[game-torsion radical theorem]\label{thm:main}
\PROVED{} In every ambient-coherent coefficient-faithful Clifford datum on
short games,
\[
 Q(t)=0,
 \qquad
 B(t,x)=B(x,t)=0
\]
for every torsion game $t$ and every short game $x$.  Consequently
\[
 Q(x+t)=Q(x)
\]
and the quadratic datum factors uniquely through
$\ShUg/\Tor(\ShUg)$.
\end{theorem}

\begin{proof}
First take a global datum.  By Theorem~\ref{thm:moews}, choose a power of two
$n$ such that $nt=0$, and choose $n$-th roots $ny=t$ and $nz=x$.  Apply
Lemma~\ref{lem:root}.  The square and anticommutator relations give
\[
 \iota(Q(t))=0,
 \qquad
 \iota(B(t,x))=0.
\]
Injectivity of $\iota$ gives the two asserted vanishings.  Symmetry of the
displayed anticommutator gives $B(x,t)=0$ as well.  Finally, the defining
polarization identity gives
\[
 Q(x+t)=Q(x)+Q(t)+B(x,t)=Q(x).
\]

For a directed family, begin in any finitely generated subgroup containing
$t$ and $x$.  Choose the roots $y,z$ in $\ShUg$ and enlarge to the finitely
generated subgroup containing $t,x,y,z$.  The preceding proof applies there.
The injective transition maps pull the vanishing back to the original
subgroup.  This is exactly where ambient coherence is used.
\end{proof}

\begin{corollary}[coefficient-independent no-go]\label{cor:coeff}
\PROVED{} Changing the commutative coefficient target cannot rescue a
game-native torsion square.  The conclusion of Theorem~\ref{thm:main} holds in
characteristic zero, characteristic two, and torsion rings such as $\Z/4$.
\end{corollary}

\begin{proof}
The proof of Lemma~\ref{lem:root} is an identity in the target algebra and
does not inspect the characteristic of $R$.
\end{proof}

\section{Consequences for the proposed escapes}

\subsection{The characteristic-two local square}

On the isolated subgroup $M=\langle *\rangle\cong\Z/2$, the declaration
$e_*^2=1$ over $\F_2$ is a valid local Clifford algebra.  Moews's theorem
supplies a short partizan game $h$ with $2h=*$.  In any coherent enlargement
containing $h$,
\[
 e_*^2=(2e_h)^2=4e_h^2=0.
\]
Thus the local square cannot survive an injective transition map.  It is a
property of the root-incomplete presentation $\langle *\rangle$, not a
quadratic invariant of the game value $*$.

The same argument eliminates the universal local record
$\operatorname{Sym}_{\F_2}(M/2M)$.  Globally,
$\ShUg=2\ShUg$, so $\ShUg/2\ShUg=0$.

\subsection{The Brown \texorpdfstring{$\Z/4$}{Z/4} route}

The older coefficient-faithfulness obstruction already shows that an odd
local value $Q(*)\in\Z/4$ collapses the scalar $2$.  Theorem~\ref{thm:main}
is stronger: even a locally faithful even value cannot extend coherently,
because the partizan roots of $*$ force $Q(*)=0$ in every coefficient ring.
Hence a genuine odd Brown phase cannot arise as the square of a globally
game-native Clifford generator.  Brown quadratic modules remain valid objects
on $\F_2$-spaces; this theorem concerns their proposed extension as faithful
Clifford square data on all short games.

\subsection{Why the Gold forms stop at the nimber core}

Finite nimber fields carry extra multiplication and trace, producing the
nonzero forms
\[
 Q_a(u)=\operatorname{Tr}(u\,u^{2^a}).
\]
These are genuinely non-tautological on the impartial field-like core.  Every
additive nimber is nevertheless a torsion short game.  Therefore
Theorem~\ref{thm:main} implies that no nonzero Gold form extends to an
ambient-coherent coefficient-faithful Clifford datum on general short games.
The obstruction is not the absence of a clever coefficient ring; adjoining
the partizan halves of nimbers kills the torsion quadratic data.

\subsection{What survives}

Quadratic data on the torsion-free quotient survives.  The implemented
mean-value and atomic-weight constructions are examples of additive game
invariants whose squares live on free directions.  The theorem predicts their
documented behavior: they are blind to torsion.

\section{Resolution and exact boundary}

\PROVED{} The \texttt{tisn} question for the current commutative-scalar
Clifford architecture has a negative answer.  A quadratic datum belonging to
short-game values must be ambient-coherent; every such coefficient-faithful
datum is Grassmann on the entire torsion subgroup and factors through the
torsion-free quotient.  This is a complete obstruction, not a finite search or
a restriction to integer coefficients.

\BOUNDARY{} The theorem deliberately does not forbid three different objects:
\begin{enumerate}[leftmargin=2em]
\item a hand-supplied table on one chosen root-incomplete subgroup;
\item the nonzero Gold forms retained inside the nimber field-like core; or
\item a directed or noncommutative game construction that is not a Clifford
      algebra over a commutative coefficient ring with additive grade one.
\end{enumerate}
The first is exactly the existing engineering API and is not game-native under
Definition~\ref{def:coherent}; the second cannot extend beyond its core; the
third changes the mathematical object and belongs to the now-resolved directed
game-semantics problem \texttt{tis}, not to a deformation of
\path{GameExterior}.  No residual coefficient target remains inside the
problem just solved.

\begin{thebibliography}{9}
\bibitem{moews}
David Moews,
\emph{The Abstract Structure of the Group of Games},
in \emph{More Games of No Chance}, MSRI Publications 42 (2002), 49--58.
\url{https://library.slmath.org/books/Book42/files/moews.pdf}.
\end{thebibliography}

\end{document}
